\documentclass[11pt]{article}
\usepackage[margin=1.1in]{geometry}
\usepackage{amsmath,amssymb,amsthm}
\usepackage{booktabs}
\usepackage{enumitem}
\usepackage{longtable}
\usepackage{array}
\usepackage{calc}
\usepackage{stmaryrd}
\usepackage{xcolor}
\usepackage{graphicx}
\usepackage{tikz}
\usetikzlibrary{arrows.meta,positioning,fit,backgrounds}
\usepackage{pgfplots}
\pgfplotsset{compat=1.18}
\usepackage[colorlinks=true,linkcolor=blue!50!black,citecolor=blue!50!black,urlcolor=blue!50!black]{hyperref}
\usepackage{natbib}

\definecolor{okgreen}{HTML}{2F7D32}
\definecolor{openamber}{HTML}{B5651D}
\definecolor{foldgray}{HTML}{4A4A4A}

\newtheorem{theorem}{Theorem}[section]
\newtheorem{lemma}[theorem]{Lemma}
\newtheorem{proposition}[theorem]{Proposition}
\newtheorem{corollary}[theorem]{Corollary}
\newtheorem{conjecture}[theorem]{Conjecture}
\theoremstyle{definition}
\newtheorem{definition}[theorem]{Definition}
\newtheorem{example}[theorem]{Example}
\newtheorem{openproblem}{Open Problem}
\theoremstyle{remark}
\newtheorem{remark}[theorem]{Remark}
\newcommand{\status}[1]{\textcolor{blue!45!black}{\textsf{[\,#1\,]}}}

\newcommand{\NN}{\mathbb{N}}
\newcommand{\CLS}{\mathrm{CL}(S)}
\newcommand{\LF}{L_{F}}
\newcommand{\LO}{\mathrm{lo}}
\newcommand{\bhs}{\mathrm{BeT}}
\newcommand{\W}{\mathrm{W}}
\newcommand{\hnf}{\mathrm{hnf}}
\DeclareMathOperator{\spine}{spine}
\DeclareMathOperator{\arity}{arity}
\DeclareMathOperator{\Drop}{Drop}


\providecommand{\tightlist}{\setlength{\itemsep}{0pt}\setlength{\par$SKI$p}{0pt}}
\setlength{\emergencystretch}{3em}
\title{A Self-Interpreter for $SKI$ Written in $SKI$:\\Extensional Semantics for New Symbols}
\author{N.\ E.\ Davis}
\date{\today}

\begin{document}
\maketitle

\begin{abstract}
A symbol added to an object language has no meaning until an interpreter
gives it one. We make that authorship visible at the smallest scale we
can find: a self-interpreter for $SKI$ combinatory logic written as a
closed term of the basis $\{S, K, I\}$, 618 atoms, which runs on an
encoding of itself in 504{,}930 contractions. Extending the object
language with a fifth symbol forces the interpreter's result type to
grow, and whether the new outcome is exposed as a value or absorbed into
divergence is decided by a single arm of the interpreter, a few atoms at
one named site, with the two variants provably agreeing on every term
that never reaches it. The same site, made a parameter, becomes a
nonlocal lookup in the style of Nock's scry: the resolver is passed as an
argument, a lookup may block, and a namespace of answers accumulates
across re-runs, keyed by a 240-atom structural equality on encoded
paths. Two structural results frame the construction. A census over the
$\{B, C, K, W\}$ basis shows that representation needs weakening, that
iteration needs contraction, and that contraction alone already costs
normalization. And the one step no $SKI$ term can perform is quotation of
a live term, by a theorem of Jay and Given-Wilson; above that boundary
everything, including metacircular evaluation, is closed.
\end{abstract}

\section{Introduction}\label{sec:intro}

Metacircular evaluation permits the separation of the evaluation
mechanism from the object language, allowing for controlled
experimentation with different evaluation strategies and the observation
of intermediate states without affecting the host environment.  An
interpreter can be employed to handle states such as partial evaluation
without interfering with the host environment, particularly in cases in
which the host environment itself may not be capable of representations
such as partially evaluated terms, intermediate representations of the
object language, crashes, and other exceptional states.

We construct a symbolic step-evaluator for the $SKI$ combinator calculus
as a closed term of the basis $\{S,K,I\}$.  This evaluator performs
symbolic evaluation by repeatedly applying the reduction rules of the
$SKI$ calculus until a normal form is reached.  Because the interpreter is
capable of symbolic representation beyond simply the basis $\{S,K,I\}$,
we can add in functionality such as handling additional combinators,
supporting custom reduction strategies, and managing exceptional states
like crashes, undefined terms, or as-yet-unknown referentially
transparent lookups.

The current line of research is motivated by the question of how
amenable the $SKI$ combinator calculus is to metacircular evaluation,
and what design patterns and implementation techniques could facilitate
such evaluation.  We do not anticipate such a line of inquiry to be
competitively efficient versus other combinator systems or higher-level
representations, but rather as a tool for further exploration of the
metacircular evaluation process and its possibilities for simple
programs.  This paper makes three contributions:
\begin{enumerate}[nosep]
\item A self-interpreter for $\{S,K,I\}$ in $\{S,K,I\}$.  This evaluator
  is a step-evaluator over encoded terms rather than a decoder, together
  with the structural census that says what each rule of the calculus is
  bought for (\S\ref{sec:interp}, \S\ref{sec:census}).
\item The observation that extending the object language forces the
  interpreter's outcome type to grow, and that the meaning of the new
  outcome is authored at one site, with a lemma that nothing else changes
  (\S\ref{sec:authored}).
\item Site as parameter.  This is a value lookup whose resolver is an
  argument, which can block, and whose namespace of answers grows across
  re-runs over structurally compared paths (\S\ref{sec:scry}).
\end{enumerate}
Section~\ref{sec:related} places the work in the metacircular tradition;
\S\ref{sec:mock} reads the construction against virtualization in the
Nock combinator calculus. The artifact and its reproduction are
described in Appendix~\ref{app:artifact}.

\section{Metacircular evaluation and self-interpretation}\label{sec:related}

While many languages nominally support it, metacircular evaluation as a
first-class strategy has only occasionally been applied outside
of the Lisp family, where it has been used to implement
features such as macro systems, partial evaluation, and sandboxed
execution environments.  Essentially, each nested interpreter acts as a
virtual machine for the next level of interpretation.  Lisp's
homoiconic nature is particularly apt for such virtualization, since the
interpreter can directly treat its interpretand as data, enabling
Lisp's powerful reflective and metaprogramming capabilities.  (In this
connexion, the collapsing towers of \citet{AminRompf2018} will be of
special interest.)

The lineage runs from \citet{McCarthy1960}, whose universal function
\texttt{eval} was written in the language it evaluates, and in
particular from the Lisp 1.5 manual \citep{Lisp15}, which specifies the
language by that function.  \citet{Reynolds1972} gave the form its
critical treatment:  a definitional interpreter written in the style of
its own object language is metacircular; the interpreted language
inherits the defining language's evaluation order unless that order is
made explicit.  (The present work meets that observation directly in
\S\ref{sec:fuel}, where the object-level laziness of $K$ is
inherited rather than implemented.)  \citet{SteeleSussman1978} and
\citet{SICP} made the metacircular evaluator a standard instrument for
experimenting with semantics, the very use to which we put it.

The Nock combinator calculus (est.\ 2002--2006) explicitly adopted
self-application as a fundamental mechanism for metacircular evaluation
from its earliest inception \citep{Nock4K}.  This mechanism was
originally intended as a ``hyper-Turing'' operator \citep{Davis2025},
in which a query for a referentially transparent value could answer
directly, with a permanent refusal, or with a deferral.

The tower as such entered with \citet{Smith1984}, whose 3-Lisp runs each
level of a reflective tower on an interpreter at the level above;
\citet{WandFriedman1988} gave the tower a non-reflective description, and
\citet{Asai1996} defined the language Black (in whose style
\citeauthor{AminRompf2018}'s Purple is cast).  \citet{AminRompf2018}
treated the tower itself as the object of study.
Their multi-level $\lambda$-calculus provides staging constructs and
stage polymorphism, so that a single evaluator acts as an
interpreter or as a compiler depending on runtime parameters, and their
metacircular evaluator Pink collapses arbitrarily many levels of
self-interpretation into one compiled program; correctness of
self-interpretation and multi-level Jones optimality are stated without
formal proof but backed by experiment (their Propositions~4.2 and~4.6).
Semantic modifications introduced at any level of the tower, such as
instrumentation, CPS conversion, or reflection in the style of Black,
survive the collapse. Their objective is performance, the removal of
interpretive overhead, and their semantics is value-producing throughout:
there is no fuel, no failure, and no outcome that has to be represented
rather than computed. The present work complements their approach. We
retain the levels with payment (\S\ref{sec:mock}), and we ask what an
interpreter does in the case where its object language gains a construct
whose meaning the interpreter must author (\S\ref{sec:authored}); this
question does not arise in a setting without a third outcome. The two
share the shape of the mechanism:  an interpreter as an ordinary program
in the calculus it interprets, with one place at which its behaviour is
substituted.  \S\ref{sec:scry} builds upon this foundation.

In the $\lambda$-calculus proper, \citet{Barendregt1991} stated the
requirement on a self-interpreter, $\mathbf{E}\,\ulcorner M\urcorner = M$,
and notes that the coding $\ulcorner\cdot\urcorner$ is a metalevel
operation rather than a term; that boundary is located by
a theorem of \citet{JayGivenWilson2011} for $SKI$ (\S\ref{sec:boundary}).
\citet{BerarducciBohm1993} show that the
choice of encoding is load-bearing, obtaining a self-interpreter that
itself has a normal form, and \citet{BrownPalsberg2016} carry
self-interpretation into System~F$_\omega$, in which typing rather than
quotation forms the barrier. Futamura's first projection
\citep{Futamura1971} handles the specialization of an interpreter to a
fixed program; \citet{JonesGomardSestoft1993} provide an standard
account of this work. \citet{Jones1996} identified the optimality
condition which must be met by a specializer to remove interpretive
overhead entirely.

Self-interpretation for a Turing-complete system is as old as
universality: \citet{Kleene1936} gave the normal form theorem from which
a universal function follows, and \citet{Turing1936} the universal
\citet{Mogensen1992} built a compact self-interpreter and
self-reducer in the untyped $\lambda$-calculus; \citet{JayPalsberg2011}
built a typed one, in a combinatory calculus with a factorisation
operator though not in $\{S,K\}$. (In particular, their distinction
between self-\emph{recognisers}, which recover a term from its
representation, and self-\emph{enactors}, which execute it, is
necessary here.) The closest prior construction for $SK$ is
\citet{Tromp2006}, whose Theorem~3 gives a $\lambda$-term $F$ of size
124 with $F\,C\,(\tilde M : N) = C\,M\,N$ for a bitstring encoding
$\tilde M$ of a combinator $M$, and whose Theorem~4 translates it into
a 263-bit self-interpreter for combinatory logic; with $C = I$ it is a
self-enactor. Architecturally $F$ is a decoder: it reconstructs live
syntax from bits and hands it to a continuation, after which reduction
proceeds under the ambient calculus's own rules. It carries no
object-level $S$ or $K$ rule in its body, and so cannot implement a step
relation that differs from the host's, which is what we carry out in
\S\ref{sec:authored}. The same paper contains a
symbolic normal-form reducer over Scott-encoded de Bruijn terms for the
$\lambda$-calculus, used only as a size constant.

Our novelty is narrow; a symbolic step-evaluator for $SKI$ in
the $\lambda$-calculus is a textbook exercise, and bracket abstraction
carries it into $\{S,K,I\}$ mechanically.  What we add is a
concrete fueled expression of 618 atoms with its step rules as data
manipulation rather than re-spliced host syntax, applied to an encoding
of itself and run to completion.  On top of this we add a census
(\S\ref{sec:census}), the one-site authorship and faithfulness lemma
(\S\ref{sec:authored}), and the substitutable, blocking, compound-path
lookup (\S\ref{sec:scry}), which we have not found elsewhere.
(\citeauthor{AminRompf2018}'s outcome space has no failure case;
\citeauthor{JayGivenWilson2011}'s result concerns the
limitations of $SK$ over live terms rather than what an interpreter
can choose to expose.)

\section{The interpreter}\label{sec:interp}

\subsection{Encoding and construction}\label{sec:construction}

The $\{S,K,I\}$ terms are the object language of the interpreter; in
the terminology of \citet{Smullyan1985} the standard definitions are:
\begin{align*}
\text{Idiot bird,}\;I & = \lambda x.x \\
\text{Kestrel,}\;K & = \lambda x.\lambda y.x \\
\text{Starling,}\;S & = \lambda x.\lambda y.\lambda z.x\,z\,(y\,z)
\end{align*}

Object terms are represented by Scott encoding
\citep{Scott1962,Jansen2013}, in which a datum of an inductive type is a
function that takes one continuation per constructor and applies the
continuation for its own constructor to its own components. The type of
object terms has four constructors, three leaves and application, so an
encoded term takes four continuations, in the fixed order $S$, $K$, $I$,
application:
\begin{align*}
\langle S\rangle &= \lambda s.\lambda k.\lambda i.\lambda z.\;s \\
\langle K\rangle &= \lambda s.\lambda k.\lambda i.\lambda z.\;k \\
\langle I\rangle &= \lambda s.\lambda k.\lambda i.\lambda z.\;i \\
\langle t\,u\rangle &= \lambda s.\lambda k.\lambda i.\lambda z.\;z\,\langle t\rangle\,\langle u\rangle
\end{align*}
A leaf discards all four continuations but its own. An application
applies the fourth continuation to the encodings of its two children;
this is how the interpreter walks a term, by applying $\langle t\rangle$
to four continuations as a case analysis on the head constructor of $t$.
For example, $\langle K\rangle$ applied to four distinct atoms
$X_1\,X_2\,X_3\,X_4$ reduces to $X_2$, and $\langle I\,K\rangle$ applied
to the same four reduces to $X_4\,\langle I\rangle\,\langle K\rangle$.
Bracket-abstracted into $\{S,K,I\}$, the encodings are closed
combinators: $\langle K\rangle$ is the five atoms $K\,(S\,(K\,K)\,K)$;
$\langle I\rangle$ is $K\,(K\,K)$; $\langle S\rangle$ is
$S\,(K\,K)\,(S\,(K\,K)\,K)$; and $\langle I\,K\rangle$ is a 40-atom term
which bears some discussion.  The bracket abstraction of
$\langle I\,K\rangle = \lambda s.\lambda k.\lambda i.\lambda z.\;z\,\langle I\rangle\,\langle K\rangle$
reads as three components composed together:
\begin{enumerate}[nosep]
  \item  The first 32 atoms are the application constructor $\text{encA}$
    itself, $\lambda t.\lambda u.\lambda s.\lambda k.\lambda i.\lambda z.\;z\,t\,u$.
  \item  Three leading prefixes $S\,(K\,(S\,(K\,K)))$ push the three
    unused binders $s$, $k$, $i$ past the body.
  \item  Finally, $z$ is applied to $t$ and $u$ in the inner term, $\lambda s k i z. k$.
\end{enumerate}
The full expression thus reads,
\begin{align*}
    & S \\
    & (K (S (K K))) (S (K (S (K K))) (S (K (S (K K))) (S (S (K S) (S (K K) (S (K S) (S (K (S I)) K)))) (K K)))) \\
    & (K (K K)) \\
    & (K (S (K K) K))\text{,}
\end{align*}
which is depicted as a tree diagram in Figure~\ref{fig:uqik}, where it is
discussed in context.  When
applied to any four terms, this expression hands $K (K K)$ and
$K (S (K K) K)$, which are $\langle I\rangle$ and $\langle K\rangle$,
to the fourth of them.

The encoding $\langle\cdot\rangle$ itself is a
metalevel operation, for the reason given in \S\ref{sec:boundary}.
Lists, the \texttt{Maybe} of the result type, and the fuel numeral are
Scott-encoded the same way:
\begin{itemize}[nosep]
  \item  $\texttt{nil} = \lambda n.\lambda c.\,n$;
  \item  $\texttt{cons}\;h\;t = \lambda n.\lambda c.\,c\,h\,t$;
  \item  $\texttt{nothing} = \lambda n.\lambda j.\,n$;
  \item  $\texttt{just}\;x = \lambda n.\lambda j.\,j\,x$;
  \item  $\texttt{zero} = \lambda z.\lambda s.\,z$; and
  \item  $\texttt{suc}\;n = \lambda z.\lambda s.\,s\,n$.
\end{itemize}
Every definition below
is written with named binders and expanded to a closed $\{S,K,I\}$ term
by bracket abstraction, with recursion through the fixpoint combinator
$Y$. The constructed interpreter, called \texttt{whnfF}, consists of a
618-atom term; its unquoter \texttt{UQ}, which maps encoded data back to
a live term, contains 43 atoms.  \texttt{whnfF} is given in full in
Appendix~\ref{app:expr} and depicted graphically in
Figure~\ref{fig:parts}.  An evaluation is illustrated in
Figure~\ref{fig:mechanism}.

\texttt{whnfF} defines a fueled interpreter that evaluates encoded terms to
weak head normal form.  The fuel loop \texttt{whnfF n $\langle t\rangle$}
peels one \texttt{suc}
from the numeral per step-attempt and returns \texttt{nothing} when the
numeral is exhausted. Each attempt calls \texttt{step}, which walks the
spine of the encoded term to its head leaf, collecting the arguments
into a list, and dispatches on the leaf: the arm for $I$ needs one
argument, for $K$ two, for $S$ three. A saturated arm performs one
contraction as data, building $\langle x\,z\,(y\,z)\rangle$ in the case
of $S$, and the rebuilder \texttt{rb} re-applies the untouched arguments;
the loop continues on the result with the remaining fuel. An unsaturated
head means the term is in weak head normal form, and the loop returns it
as \texttt{just $\langle t\rangle$}. This is normal-order head reduction
with one contraction per attempt, which matches the host reducer's count
exactly.

\begin{figure}[tb]
\centering
\includegraphics[width=\linewidth]{build/fig/fig-engine-parts.pdf}
\caption{The interpreter's components as terms: the
spine walker \texttt{sp} (126 atoms), the rebuilder \texttt{rb} (74), the
three step arms \texttt{stepI} (92), \texttt{stepK} (105), \texttt{stepS}
(237), and the unquoter \texttt{UQ} (43). Each panel is scaled to its
frame, so sizes are not comparable across panels. Function edges are
green and argument edges blue; $S$ is a filled circle, $K$ a square, $I$
a ring. The $S$ arm is substantially larger because it builds
$\langle x\,z\,(y\,z)\rangle$ as data.}
\label{fig:parts}
\end{figure}

\begin{figure}[tb]
\centering
\includegraphics[width=\linewidth]{build/fig/fig-engine-mechanism.pdf}
\caption{The \texttt{whnfF} evaluation process. The fuel loop peels one
\texttt{suc} per step-attempt; \texttt{step} walks the spine to the head
leaf, collecting the arguments; dispatch selects the arm for that leaf; a
saturated arm contracts and \texttt{rb} re-applies the untouched
arguments, returning \texttt{just $\langle t'\rangle$} to the loop; an
unsaturated head returns \texttt{just $\langle t\rangle$}, the value;
exhausted fuel returns \texttt{nothing}. The dashed fifth arm optionally
contains an authored site (\S\ref{sec:authored}).}
\label{fig:mechanism}
\end{figure}

\subsection{Self-application}\label{sec:tower}

The \texttt{whnfF} interpreter is applied to encoded terms to
evaluate them to weak head normal form.  This emulated process is not
efficient: it incurs significant overhead due to the emulation of the
host reduction steps.  With
fuel 3 on $\langle I\,K\rangle$ it reaches $\langle K\rangle$ in 340 host
contractions; the decoded result equals the host's own reduction of
$I\,K$. When it is run over the unquoter applied to an encoded term, it
takes 91,556 steps. Run over an encoding of itself running on $\langle I\,K\rangle$,
as a tower, it reaches weak head normal form in 504,930 contractions
with the correct value. Empirically for this case, layer one costs 340
contractionsto emulate one, and layer two costs about 1,485 times layer
one, the ordinary per-level overhead of naive self-interpretation.
Figure~\ref{fig:motion} shows the first run in progress and
Figure~\ref{fig:whnff} the interpreter as a term.

\begin{figure}[b]
\centering
\includegraphics[width=\linewidth]{build/fig/fig-engine-motion.pdf}
\caption{The interpreter running \texttt{whnfF 3 $\langle$I K$\rangle$},
after 0, 20, 60, 120, 200, 300 and 340 host contractions, in the spiral
view. The spine curls as $Y$ unfolds and the encoded term is walked; the
tree reaches about 111{,}000 atoms at step 300 and collapses to the
answer, \texttt{just $\langle$K$\rangle$}, at step 340.}
\label{fig:motion}
\end{figure}

\begin{figure}[t]
\centering
\includegraphics[width=\linewidth]{build/fig/fig-pair-whnff.pdf}
\caption{The interpreter \texttt{whnfF}: the angle is in-order
leaf position, the radius is depth (which bottoms out at 48 levels).
Each fern-like arm contains one deeply nested argument as produced by
the bracket abstraction. Appendix~\ref{app:expr} gives the term in full.}
\label{fig:whnff}
\end{figure}

\subsection{The encoding boundary}\label{sec:boundary}

$S\,K\,K$ and $I$ agree on every argument, yet they are distinct normal
forms as terms of weak combinatory logic. A quotation function from live
terms to encoded data must be able to distinguish them, which cannot be
done by any pure $SKI$ term.

\begin{proposition}\label{prop:quote}
No closed term $Q$ over $\{S,K,I\}$ satisfies $Q\,M \twoheadrightarrow
\langle M\rangle$ for every closed $M$.
\end{proposition}

\begin{proof}
\citet[Thm.~3.2]{JayGivenWilson2011} prove that factorisation of
$SK$-combinators is a symbolic computation not representable by any
$SK$-combinator; the proof is stated for $\{S,K\}$ and goes through
unchanged with $I$ primitive. Quotation would give factorisation: $Q$
together with the unquoter \texttt{UQ} of \S\ref{sec:construction}
decomposes any term into its head and arguments by Scott dispatch on
$\langle M\rangle$. Directly: under the translation to the
$\lambda$-calculus, $\llbracket Q\rrbracket(\lambda x.x)$ is a single
term, yet $Q\,(S\,K\,X)$ would yield $\langle S\,K\,X\rangle$, from which
$X$ is extractable, so every $X$ would share a reduct, contradicting
that $S$ and $K$ are distinct normal forms.
\end{proof}

Thus the encoder is a (better, the only) host function.  Quotation is
a one-time cost which must be paid to enter the representation.
Everything on the other side of the boundary, the unquoter, the
interpreter, their self-application, and the lookup machinery of
\S\ref{sec:scry} including equality on encoded paths, is ordinary
$\{S,K,I\}$. Figure~\ref{fig:uqik} shows the decoder beside a datum:
in this calculus they are the same kind of thing.

\begin{figure}[htbp]
\centering
\includegraphics[width=\linewidth]{build/fig/fig-pair-uq-ik.pdf}
\caption{The decoder and a datum. \texttt{UQ}, 43 atoms, and the Scott
encoding $\langle I\,K\rangle$ of the two-atom term $I\,K$, 40 atoms, each
in spiral and every-atom views. A datum and a program look alike because
they are the same kind of thing; only the encoding step that produced
$\langle I\,K\rangle$ lies outside the calculus.}
\label{fig:uqik}
\end{figure}

\section{The structural cost}\label{sec:census}

What must any self-interpreter for $SKI$ be able to do, structurally?
$S$ bundles two capabilities into one rule, composition and the
duplication of its third argument, and never discards; $K$ supplies
discarding.  The two capabilities that matter here, discarding and
copying, coincide with the structural rules of weakening and contraction
in substructural logic.
By utilizing the basis $\{B, C, K, W\}$ for a comparative rewrite, we
access the capabilities directly: $B$ composes, $C$ swaps, $K$ discards
an argument, and $W$ duplicates an argument.  We are interested in
analyzing the structural cost of each component of the interpreter and
the representation layer.  Hence by rebuilding each component
over the $\{B, C, K, W\}$ basis, we see which structural rule it
requires, whether weakening ($K$) or contraction ($W$). We also count,
independently of any abstraction algorithm, how often each binder occurs
in each component's definition: a component needs weakening iff some
binder occurs zero times, and contraction iff some binder occurs more
than once. This occurrence census is the invariant measurement; atom
counts are convention-relative (with $I$ primitive the interpreter core
shows seven $W$ atoms, with $I = W\,K$ ten).

\paragraph{The representation layer is affine.} All ten Scott
constructors, the four term constructors together with \texttt{nil},
\texttt{cons}, \texttt{nothing}, \texttt{just}, \texttt{zero}, and
\texttt{suc}, have no binder occurring twice: a constructor selects one
continuation and discards the rest, which is weakening by construction.
The representation lives in $BCK$.

\paragraph{The interpreter core needs contraction at five of eight
components.} There are seven surplus occurrences in all: the unquoter's
$f$ (twice, in $(f\,t)(f\,u)$), the spine generator's accumulator (four
times), the rebuild generator's head (twice), the $S$-arm's $z$ (twice,
in $x\,z\,(y\,z)$, the $S$ rule itself reappearing one level up), and the
fuel loop's term (twice). Any $BCKW$ abstraction must perform at least
that many contractions.

\paragraph{Each rule is needed, and one of them costs normalization.}
Removing weakening removes dispatch: a Scott constructor is an $n$-ary
discard. Removing contraction removes iteration: $Y$ duplicates its
argument, and the affine fragment $BCK$ is strongly normalizing, so it
has no fixpoint combinator and no fuel loop. (The fuel itself, a Scott
numeral, is affine; iteration comes from $Y$.) The two are not
symmetric. Contraction alone already admits divergence $\bot$:
$\Omega = (S\,I\,I)(S\,I\,I)$ contains no $K$, and in $BCKW$ the term
$W\,W\,W$ reduces to itself with no weakening anywhere, so the relevant
fragment $BCIW$ is not normalizing while $BCK$ is. Weakening is bought
for representation and costs nothing in normalization; contraction is
bought for iteration and is exactly what makes non-termination
possible. Hence any self-interpreter for $SKI$ must be written in a
fragment where non-termination exists, and its outcome space cannot be
$\{\text{value}\}$: it is at least $\{\text{value}, \bot\}$,
which the next section makes observable by introducing fuel. The third
outcome is not a design choice of this artifact but a consequence of
self-interpretation.

\paragraph{The census also explains the artifact's shape.} The
representation is cheap because data only selects: every constructor and
every case analysis is weakening, and the whole encoding of a term is
built from discards. The cost is concentrated in five places, and all
five are either recursion (the unquoter, the spine walker, the
rebuilder, and the fuel loop each pass themselves or their state on
twice) or the $S$ rule reappearing one level up inside the arm that
implements it, which is why that arm is the largest component
(Figure~\ref{fig:parts}). Two practical consequences follow: the
representation layer can be carried unchanged into an affine calculus,
with only the loop needing the full basis; and any effort to shrink the
interpreter should be spent on those five sites, since everything else
already bears the cost of a selection.

\section{Extensional authored semantics}\label{sec:authored}

In a metacircular evaluator augmenting interpretation,
additional outcomes can be observed and handled explicitly, rather than
being hidden as divergence.  This permits concepts like fuel and
\texttt{Maybe} for partial evaluation to be introduced to a basis
that without these concepts natively.

\subsection{Fuel and the third outcome}\label{sec:fuel}

An unfueled interpreter's outcome space is $\{\text{value},
\bot\}$, and divergence cannot be observed from inside a
finite run. Fuel converts it to $\{\text{value}, \text{timeout}\}$,
which are both observable. Such limitation is local and carries a clear
price: at fixed fuel $k$, one observation does not separate a divergent
term from one that needs more than $k$ steps. $\Omega$ times out at
every $k$; $I^{31}K$ times out for $k \le 31$ and yields a value from
$k = 32$, since fuel $k$ permits $k$ step-attempts of which the last
must be the no-redex check. Convergence is semi-decidable by raising
$k$ while divergence is not (the halting problem). Object-level laziness
is inherited, so $KI\Omega$ discards $\Omega$ unevaluated under the
interpreter exactly as under host reduction because the object-level $K$
rule fires before the interpreter recurses into the argument.  (This
follows \citeauthor{Reynolds1972}'s observation.)

\subsection{The site}\label{sec:site}

Extend the object language with a fifth constructor \texttt{Err}, with
no object-level contraction rule. Its encoding is unremarkable, five
continuations instead of four. When the spine walker finds \texttt{Err}
in head position, the interpreter decides how to dispatch:
\begin{verbatim}
D("errArmAbs", ["acc"], a("errd"))  # observable: crash becomes a value
D("errArmOmg", ["acc"], OMEGA)      # hidden: crash becomes host divergence
\end{verbatim}
These arms are written $K\,(K\,(K\,I))$ and $K\,(S\,I\,I\,(S\,I\,I))$;
our two full interpreters differ at this juncture. The result $\Omega =
(S\,I\,I)(S\,I\,I)$ was selected as one of five divergent six-atom terms
in this basis.  The results of the two interpreters evaluating five
different terms are summarized in Table~\ref{tab:behavior}.

% \begin{figure}[tb]
% \centering
% \includegraphics[width=\linewidth]{build/fig/fig-arms.pdf}
% \caption{The two arms of the authored site. $K\,(K\,(K\,I))$, four atoms,
% exposes the crash as a value; $K\,(S\,I\,I\,(S\,I\,I))$, seven atoms,
% absorbs it into divergence. The two interpreters differ by exactly these
% atoms.}
% \label{fig:arms}
% \end{figure}

\begin{table}[tb]
\centering
\caption{Behavior of the two interpreters on selected object terms.
The outcome ``host diverges'' refers to the step cap of the host reducer
being reached, not a detected non-termination generally; the arm is
$\Omega$.}
\label{tab:behavior}
\begin{tabular}{lll}
\toprule
object term & exposing (\texttt{errArmAbs}) & hiding (\texttt{errArmOmg}) \\
\midrule
\texttt{Err}     & ERR  & host diverges \\
\texttt{Err K}   & ERR  & host diverges \\
\texttt{K I Err} & VAL  & VAL \\
\texttt{I K}     & VAL  & VAL \\
$\Omega$         & TIME & TIME \\
\bottomrule
\end{tabular}
\end{table}

\begin{lemma}[Faithfulness]\label{lem:faithful}
For every object term whose reduction never brings \emph{\texttt{Err}} into
head position, the exposing and hiding interpreters produce the same
result.
\end{lemma}

\begin{proof}
The two definitions differ only in the arm \texttt{errArm}, which is
referenced only from the \texttt{Err} case of the step function, which
the spine walker selects only when the head leaf is \texttt{Err}. On any
other input the two are the same program.
\end{proof}

The \texttt{K I Err} row of Table~\ref{tab:behavior} refers to an
instance in which \texttt{Err} is present but discarded by object-level
laziness before it is ever in head position. Here it highlights that the
addition of a constructor does not by itself introduce a new meaning.
The meaning is authored downstream by the author of the interpreter.

\section{Site as a parameter}\label{sec:scry}

The substitution as written in \S\ref{sec:site} affords a choice between
two source variants. A lookup in the style of Nock's ``scry'' or peek
into a bound referentially transparent namespace takes the resolver as
an argument instead. We build that in three steps, each extending the
outcome type once more.

\subsection{A synchronous resolver}\label{sec:sync}

The object language gains a fifth leaf \texttt{Scry}, of arity one,
in a five-symbol encoding parallel to \texttt{Err}. The fueled loop
gains a parameter: \texttt{wfQ E n $\langle t\rangle$}, where
$E$ is an ordinary encoded partial function from paths to values, a
\texttt{Maybe}. When the spine walker reaches \texttt{Scry p} it applies
$E$ to $p$, the path exactly as written and unreduced. A hit splices the
answer in and reduction continues. A miss is kept as a third step outcome
distinct from ``stuck'', reusing the three-way type of
\S\ref{sec:site}, rather than folded into the existing \texttt{Maybe},
which would conflate ``the resolver had no answer'' with ``no redex''.
The arm itself, in the same notation as the \texttt{Err} arms of
\S\ref{sec:site}, is
\begin{verbatim}
D("scHit",    ["rest","v"],     A(a("stepped"), A(a("rbQ"), v("v"), v("rest"))))
D("stepScQ1", ["e","p","rest"], A(v("e"), v("p"),
    a("errd"),                      # nothing    -> the miss
    A(a("scHit"), v("rest"))))      # just v     -> splice and continue
\end{verbatim}
that is, apply the resolver to the path and hand it two continuations:
on \texttt{nothing}, the miss; on \texttt{just v}, splice \texttt{v} in
and rebuild. The parameter is the resolver $E$; when $E$ is fixed at
compile time to one answer, the two arms become the \texttt{Err} arms of
\S\ref{sec:site}.

The loop matches the base interpreter on every term without
\texttt{Scry} against three resolvers (one that answers every path by
echoing it, one that answers none, one that dispatches on the path's head
leaf and answers only $K$); hits and misses where expected; reports a
bare \texttt{Scry} as stuck rather than as either; reduces
\texttt{K K (Scry K)} to $K$ without consulting the resolver, so a site
inside a branch the program never forces is never consulted; and times
out on $\Omega$ before any site is reached. Substitutability is not
eager; since the path is unreduced, \texttt{Scry (I K)} and
\texttt{Scry K} are different paths, a point to which \S\ref{sec:paths}
recurs. A path is not a sequence and not an address: it is whatever
term occupies \texttt{Scry}'s argument position, delivered to the
resolver as an encoded term, and it bears no meaning to the interpreter
itself. Its interpretation is entirely the resolver's business. (Nock's
scry paths are defined as lists of atoms; nothing here depends on that
shape.)

\subsection{A blocking resolver}\label{sec:block}

A resolver that always answers is not a scry because a real lookup can
fail to answer \emph{yet} for a given path (such as a network call or
result of a computation). The resolver's type grows from two-way to
three-way (\texttt{oJust}, \texttt{oNothing}, \texttt{oNotYet}) and the
interpreter's outcome types grow to match, adding \texttt{pending} at the
step level and \texttt{block} at the result level. The arm gains one
continuation:
\begin{verbatim}
D("stepScN1", ["e","p","rest"], A(v("e"), v("p"),
    A(a("scHitN"), v("rest")),      # oJust v    -> splice and continue
    a("errdN"),                     # oNothing   -> the miss
    A(a("pendingN"), v("p"))))      # oNotYet    -> block, carrying the path
\end{verbatim}
The difference against the synchronous arm above is the third line.
This extension triggers the blocking phenomenon of \S\ref{sec:authored},
providing one continuation at one site.

Resuming a block captures no continuation. Because reduction is pure and
deterministic, the driver re-runs the whole term from the top under a
resolver that now also answers the missing fact. Two conditions make
that sound. The re-run's steps strictly before the block are identical
to the first run's only if the new resolver agrees with the old one on
every path the old one answered, which holds because the namespace is
append-only. Re-running is not fuel-neutral: each round replays the
previous one and then takes at least one more step, so with a fixed
budget a block can become a timeout on the next round. Recomputation is
charged for. The interpreter and each round's resolver are closed
$\{S,K,I\}$ terms; the host side provides the fact table standing in for
an external source of truth, the resume loop, and the detection of a
path with no fact.  (This starts to gesture towards a runtime handler
which facilitates $SKI$ evaluation in the context of blockable
computations.)

The namespace is accumulative. With an empty namespace and facts $S
\mapsto \texttt{Scry}\ K$, $K \mapsto I$, the term \texttt{Scry S}
blocks on $S$, learns it, re-runs, and the newly substituted
\texttt{Scry K} blocks on $K$, a fact forced by the result of resolving
the first, before the third round reaches the value $I$. A path with no
fact blocks into a distinguishable stuck state, and fuel exhaustion
remains distinguishable from a block throughout.

\subsection{Compound paths}\label{sec:paths}

The resolver of \S\ref{sec:block} dispatches on a path's head leaf only
and cannot tell $S\,K$ from $S\,K\,K$. Proposition~\ref{prop:quote} does
not force this, as that theorem concerns factorising \emph{live}
combinators. The path a resolver receives is Scott-encoded \emph{data},
a tree with five leaf tags, and equality on such data is an ordinary
structural recursion. \citeauthor{JayGivenWilson2011} said as much:
once a term is encoded, factorisation ``is a routine list operation''.
However, $SK$ cannot produce the encoding from a live term.

We build \texttt{EQ5}, a structural equality test over the
five-constructor encoding (Figure~\ref{fig:eq5}); and a resolver $p$ that
compares the incoming path against each known path in turn,
\[
p \;\mapsto\; \texttt{EQ5}\ p\ \langle k_1\rangle\ (\texttt{oJust}\
\langle a_1\rangle)\ (\texttt{EQ5}\ p\ \langle k_2\rangle\ (\texttt{oJust}\
\langle a_2\rangle)\ (\cdots\ \texttt{oNotYet})),
\]
which runs inside the blocking loop unchanged. This distinguishes a path
from its prefix, resolving \texttt{Scry (S K K)} and \texttt{Scry (S
K)} to different values under the same facts; it chains through a
compound answer, \texttt{Scry (K I)} resolving to \texttt{Scry (S K)} and
then to $I$ in three rounds; and it refuses a false hit, blocking on
\texttt{Scry (S K)} when the only fact is for $K\,S$.

\begin{figure}[htbp]
\centering
\includegraphics[width=\linewidth]{build/fig/fig-pair-eq5.pdf}
\caption{\texttt{EQ5}, structural equality on encoded paths, 240 atoms:
92 $S$, 118 $K$, 30 $I$. Spiral view, 68 atoms visible, beside the
every-atom disc.}
\label{fig:eq5}
\end{figure}

Paths are intensional, syntactic objects, compared by structural
equality of their encodings, exactly as a Nock scry path is a noun
compared by noun equality. \texttt{Scry (I K)} blocks when the only fact
is for $K$, and resolves when a fact is given for $I\,K$. Matching paths
up to conversion constitutes the known-undecidable word problem for
$SKI$; normalizing the path under fuel before lookup would be a fueled
approximation thereto and would reintroduce the third outcome inside
path resolution.

\section{Discussion: virtualization and \texttt{+mock}}\label{sec:mock}

Nock's virtualized interpreter \texttt{+mock} is interesting not because
Nock can interpret Nock; that follows once a formula computes
\texttt{nock}, since a formula already is a noun. Its interest lies
instead in that a formula can be run through a \emph{substituted}
interpreter, one that resolves a scry differently or catches a crash
and returns it as a value, while agreement with the reference
interpreter holds wherever the substitution is not exercised, and the
substitution is confined to a named site. Read that way,
\S\ref{sec:authored} is the $SKI$ analogue of Nock virtualization
(rather than the tower of \S\ref{sec:tower}). The tower demonstrates
self-interpretation and fidelity, and a general version of that
demonstration exists for value-producing towers
\citep[Props.~4.2 and~4.6]{AminRompf2018}. The two arms demonstrate
the same thing \texttt{+mock} is for: the same program, the same
interpreter apart from one arm, different observable outcomes in the
case where the outcome categories differ. Lemma~\ref{lem:faithful}
asserts/requires faithfulness of the metacircular interpreter.

The analogy with Nock breaks at quotation. \texttt{+mock} never needs
a quoting step, because a Nock formula already is the noun handed to
the alternate interpreter. $SKI$ is not homoiconic: $S$, $K$, $I$ are
opaque, and Proposition~\ref{prop:quote} states that no $SKI$ term
can turn an arbitrary $SKI$ term into data about itself. $SKI$-in-$SKI$
is therefore not symmetric the way Nock-on-Nock is. There is one
unavoidable non-$SKI$ quoting step, and virtualiation is native only
above that boundary and closed under further quoting.

Depth is bounded for the ordinary reason. Two levels already cost
504{,}930 contractions against a trivial base term, at roughly $10^3$
per level; a third level was not attempted in this research program.
(Nock-on-Nock virtualization has the same character, which is why
runtime-matched accelerated code exists for Nock as a practical
instruction set architecture.)

\section{Conclusion and open problems}\label{sec:open}

An interpreter that is an ordinary term of the calculus it interprets
makes the authorship of meaning visible: a new symbol's semantics is a
few atoms at one site, provably isolated from everything else; that
same site made into a parameter furnishes a namespace. Our construction
is small enough to be read whole (Appendix~\ref{app:expr}) and cheap
enough to extend, and several open research questions follow therefrom:

\begin{itemize}[nosep]
\item \emph{Partial evaluation.} A resolver that treats an unanswered
  path as an inert unknown, rather than as a block, turns the blocking
  loop into a partial evaluator: run over an encoding of the interpreter
  applied to a fixed program with the resolver unknown, it returns the
  program with every dispatch before the first lookup executed away,
  Futamura's first projection in its weakest form. The static prefix is
  saved exactly; what remains after the first dynamic dependency is the
  interpreter's own continuation, and removing it requires unfolding
  under an unknown head, which is a separate project.
\item \emph{Quotation above level 0.} Proposition~\ref{prop:quote} bars
  quotation of live terms only. Re-encoding an already-encoded term,
  $\langle t\rangle \mapsto \langle\langle t\rangle\rangle$, is a data
  transformation and definable, and with it a tower can be deepened from
  inside; whether this recovers the self-quotation of
  \citet{Jay2016,Jay2017} in the intensional setting is open.
\item \emph{An unquoter for the extended encoding,} and a pure-$SKI$
  driver for the blocking loop, would move the last host-side pieces of
  \S\ref{sec:block} into the calculus. The resolver's namespace is
  host-side only because its facts are compiled into a closure; an
  explicit core convention, a Scott pair of battery and payload in the
  manner of a Nock core, would make a fact an ordinary \texttt{cons} and
  the whole loop one closed term. That convention, and a subject-oriented
  surface language over it whose compile-time/runtime line for quotation
  is fixed by Proposition~\ref{prop:quote}, are the subject of separate
  work. The semantics of intensional recursion \citep{Kavvos2017} is the
  natural framework in which to state what such a driver would be.
\item \emph{Typed languages atop combinators.} The encoding of typed
  languages on top of combinatory logic, particularly in the presence
  of intensional features, can permit new kinds of type safety for
  second- and higher-order values.  The Hoon language for the Nock
  combinator calculus exemplifies this but the space has not been
  extensively explored.
\end{itemize}
While $SKI$-on-$SKI$ metacircular interpretation is hardly efficient,
the development of a toolkit of reusable components and abstractions for
$SKI$-on-$SKI$ interpretation affords a promising structure for future
exploration of combinator-based computing.

\bibliographystyle{plainnat}
\bibliography{mss}

\appendix

\section{The artifact}\label{app:artifact}

All results were produced by the scripts in the accompanying repository
against \texttt{aviary-kernel} 0.1.0, and every script checks its own
claims: each printed result is compared against a fresh-marker probe of
the reduced value rather than a printed outcome tag, and each script
exits with a nonzero status if any check fails.

\paragraph{Files.} \texttt{tower\_harness.py} defines the interpreter
and the two \texttt{Err} variants; \texttt{census.py} and
\texttt{bckw\_census.py} perform the census of \S\ref{sec:census};
\texttt{scry\_harness.py}, \texttt{scry\_namespace.py}, and
\texttt{scry\_paths.py} build the three loops of \S\ref{sec:scry}, each
loading the previous file's environment and modifying nothing in it;
\texttt{figures.py} draws every figure from the expanded terms;
\texttt{review/} holds the brute-force and construction checks cited in
the text. \texttt{ski-in-ski.ipynb} executes end to end and reproduces
the 618, 43, and 504{,}930 figures in its stored outputs;
\texttt{ski-in-ski-aviary.ipynb} builds the interpreter natively in the
\texttt{aviary} kernel as a walkthrough.

\paragraph{Sizes.} Tree-atom counts, from \texttt{aviary\_kernel.terms.size}:
\begin{center}
\begin{tabular}{lr}
\toprule
term & atoms \\
\midrule
\texttt{whnfF}, the interpreter & 618 \\
\texttt{UQ}, the unquoter & 43 \\
\texttt{wf5Abs} / \texttt{wf5Omg}, the two \texttt{Err} variants & 768 / 771 \\
\texttt{wfQ}, synchronous resolver loop & 950 \\
\texttt{wfN}, blocking resolver loop & 1{,}066 \\
\texttt{EQ5}, equality on encoded paths & 240 \\
\bottomrule
\end{tabular}
\end{center}
The interpreter script's banner prints 348 and 62 for the first two;
those are counts of DAG nodes, atoms plus application nodes deduplicated
by identity, which is why the second exceeds 43.

\paragraph{Targets.} \texttt{python3 tower\_harness.py --target t0}
(340 contractions), \texttt{t1} (91{,}556), \texttt{t2 --kk 500}
(504{,}930; any object fuel of at least 341 gives this count, and the
script refuses less), and \texttt{t3} (the table of \S\ref{sec:site}).
The three scry scripts take no arguments. The enumeration of divergent
terms is \texttt{review/brute2.py}; the six-atom arm is
\texttt{review/omg6.py}.

\section{The interpreter as an expression}\label{app:expr}

The terms of Figures~\ref{fig:whnff} and~\ref{fig:uqik}, written out.
Application is juxtaposition and associates to the left; parentheses
appear only where association is to the right; line breaks fall at
spaces and carry no meaning. Both are exactly what bracket abstraction
produces from the named definitions, so their atom counts are
algorithm-relative in the sense of \S\ref{sec:census}. Both begin with
the same fourteen atoms, \texttt{S (K (S I I)) (S (S (K S) K) (K (S I
I)))}, which is $Y$: each is a fixpoint of its generator.

\texttt{UQ}, the unquoter, 43 atoms:
\begin{verbatim}
S (K (S I I)) (S (S (K S) K) (K (S I I))) (S (K (S (S (S (S I (K S)) (K K))
(K I)))) (S (K K) (S (S (K S) (S (K (S (K S))) (S (K K)))) K)))
\end{verbatim}

\texttt{whnfF}, the interpreter, 618 atoms:
\begin{verbatim}
S (K (S I I)) (S (S (K S) K) (K (S I I))) (S (K (S (S (K S) (S (K K) (S I (K
K)))))) (S (K K) (S (K (S (S (K S) (S (K K) (S (S (S (S (S (S (K (S I I)) (S
(S (K S) K) (K (S I I))) (S (K (S (S (K S) (S (S (K S) (S (S (K S) (S (S (K
S) K) (K (S (K (S (K K))) (S (K (S (K K))) (S (K (S I)) K)))))) (K (S (K K)
(S (K (S (K K))) (S (K (S I)) K)))))) (K (S (K K) (S (K K) (S (K (S I))
K)))))))) (S (K K) (S (S (K S) (S (K K) (S (K S) (S (K (S (K S))) (S (K
K)))))) (K (S (K K) (S (K (S (S (K (S (K K))) (S (S (K S) (S (K K) (S (K S)
(S (K (S I)) K)))) (K K))))) K))))))) (K K)) (K (S (S I (K K)) (K (S (K (S
(S I (K K)))) (S (K K) (S (K (S (K (S (S I (K K)))))) (S (K (S (K K))) (S (K
(S (K (S (K (S (K (S (K K) (S (K (S I)) K))))))))) (S (K (S (K (S (K (S (K
(S I I)) (S (S (K S) K) (K (S I I))) (S (K (S (S (K S) (S (K (S I)) K)))) (S
(K (S (K K))) (S (S (K S) (S (K K) (S (K S) K))) (K (S (K (S (K K))) (S (K
(S (K K))) (S (K (S (K K))) (S (S (K S) (S (K K) (S (K S) (S (K (S I)) K))))
(K K))))))))))))))) (S (S (K S) (S (K K) (S (K S) (S (K (S (K (S (K (S (K
K))) (S (K (S (K K))) (S (K (S (K K))) (S (S (K S) (S (K K) (S (K S) (S (K
(S I)) K)))) (K K)))))))) (S (K (S (K K))) (S (K (S (K K))) (S (K (S (K K)))
(S (S (K S) (S (K K) (S (K S) (S (K (S I)) K)))) (K K))))))))) (K (S (K (S
(K K))) (S (K (S (K K))) (S (K (S (K K))) (S (S (K S) (S (K K) (S (K S) (S
(K (S I)) K)))) (K K))))))))))))))))) (K (S (S I (K K)) (K (S (K (S (S I (K
K)))) (S (K K) (S (K K) (S (K (S (K (S (K K) (S (K (S I)) K))))) (S (K (S I
I)) (S (S (K S) K) (K (S I I))) (S (K (S (S (K S) (S (K (S I)) K)))) (S (K
(S (K K))) (S (S (K S) (S (K K) (S (K S) K))) (K (S (K (S (K K))) (S (K (S
(K K))) (S (K (S (K K))) (S (S (K S) (S (K K) (S (K S) (S (K (S I)) K)))) (K
K)))))))))))))))))) (K (S (S I (K K)) (K (S (K (S (K (S (K K) (S (K (S I))
K))))) (S (K (S I I)) (S (S (K S) K) (K (S I I))) (S (K (S (S (K S) (S (K (S
I)) K)))) (S (K (S (K K))) (S (S (K S) (S (K K) (S (K S) K))) (K (S (K (S (K
K))) (S (K (S (K K))) (S (K (S (K K))) (S (S (K S) (S (K K) (S (K S) (S (K
(S I)) K)))) (K K))))))))))))))) (S (K K) (S (K (S I)) K))))))) K)))
\end{verbatim}

\end{document}
